
\documentclass{article}
\usepackage{colortbl}
\usepackage{makecell}
\usepackage{multirow}
\usepackage{supertabular}

\begin{document}

\newcounter{utterance}

\centering \large Interaction Transcript for game `codenames', experiment `association\_difficult', episode 4 with reflexion{-}Qwen3.5{-}27B.
\vspace{24pt}

{ \footnotesize  \setcounter{utterance}{1}
\setlength{\tabcolsep}{0pt}
\begin{supertabular}{c@{$\;$}|p{.15\linewidth}@{}p{.15\linewidth}p{.15\linewidth}p{.15\linewidth}p{.15\linewidth}p{.15\linewidth}}
    \# & $\;$A & \multicolumn{4}{c}{Game Master} & $\;\:$B\\
    \hline

    \theutterance \stepcounter{utterance}  
    & & & \multicolumn{2}{p{0.3\linewidth}}{
        \cellcolor[rgb]{0.9,0.9,0.9}{
            \makecell[{{p{\linewidth}}}]{
                \texttt{\tiny{[GM$|$GM]}}
                {'hidden': {'team': ['polyester', 'gymnastics', 'france', 'birch', 'hat'], 'innocent': ['ostrich', 'volleyball', 'korea', 'pants', 'archbishop', 'acid'], 'opponent': [], 'assassin': ['coal']}, 'revealed': {'team': {'team': [], 'innocent': [], 'opponent': [], 'assassin': []}, 'opponent': {'team': [], 'innocent': [], 'opponent': [], 'assassin': []}}}
            }
        }
    }
    & & \\ \\

    \theutterance \stepcounter{utterance}  
    & & \multicolumn{4}{p{0.6\linewidth}}{
        \cellcolor[rgb]{0.9,0.9,0.9}{
            \makecell[{{p{\linewidth}}}]{
                \texttt{\tiny{[P1$\langle$GM]}}
                \texttt{You have attempted this same task before and did not succeed. These are the plans you wrote after those attempts.} \\
\texttt{Trial \#1: New plan: The previous attempt was aborted because the response included extra text ("You: ") before the required format, violating the instruction to "never include any other text than is required in this form." Additionally, the clue "tree" was too narrow, only covering two words. Next time, I will output strictly the two lines starting with "CLUE:" and "TARGETS:" with no prefix or suffix. I will also select a broader clue, such as "fabric" or "material," to cover more target words like polyester, wool (opponent, so avoid), hat, and potentially others, ensuring the clue is semantically related but not a substring of the targets. Specifically, I will use "fabric" to target polyester and hat, or "bird" for swan, ensuring the format is exact:} \\
\texttt{CLUE: <WORD>} \\
\texttt{TARGETS: <WORD1>, <WORD2>, ...} \\
\\ 
\texttt{Trial \#2: New plan: My previous strategy failed because I provided clues that were too broad or ambiguous, leading my partner to guess opponent words (e.g., "fabric" triggered "wool", "sport" triggered "hockey", "tree" triggered "pine"). I also wasted turns by targeting only one word at a time when multiple words could be safely grouped. Next time, I will prioritize high{-}specificity clues that uniquely identify my remaining words while explicitly avoiding associations with opponent, distractor, or kill words. I will group multiple safe targets into a single clue to maximize efficiency. For example, instead of "sport" (which risks "hockey"), I will use "gym" or "rings" for "gymnastics". Instead of "tree" (which risks "pine"), I will use "wood" or "timber" for "cedar" and "birch" only if "pine" is no longer a distractor, or target them individually with specific types like "evergreen" if safe. I will strictly verify that the clue does not semantically link to any remaining opponent words before submitting.} \\
\\ 
\texttt{Trial \#3: New plan: My previous attempt failed because I only targeted two words ("polyester", "hat") with the clue "fabric", leaving "gymnastics", "france", and "birch" unguessed, which likely allowed the opponent to win first. I need to maximize the number of targets per clue while ensuring safety. I will group "gymnastics" and "france" under the clue "olympics" (as France hosted the 2024 Olympics and gymnastics is a key event), and "birch" and "polyester" under "material" or "tree" (if safe from distractors like "ostrich" or "acid"). Specifically, I will start with a clue like "olympics" targeting "gymnastics" and "france" to clear two words immediately, then follow up with a clue for the remaining words, ensuring I never leave words unaddressed if a safe multi{-}word clue exists. I must verify that "olympics" does not trigger "korea" (opponent) or "volleyball" (distractor) before using it. If "olympics" is risky, I will use "sport" for "gymnastics" and "country" for "france" separately, but prioritize grouping whenever possible to speed up the game.} \\
\\ 
\\ 
\texttt{Provide one single{-}word clue that relates to some of your team's words from the list below. You can choose to how many and to which words your clue relates to.} \\
\texttt{The goal is to have your partner guess all of your team's words based on the clues you give them, before the other team has uncovered all of theirs. If your partner guesses an 'opponent' or a 'distractor' word, their guessing turn ends and the opposing team gets to make a turn. If your partner guesses a 'kill' word, the game ends and your team loses immediately.} \\
\texttt{The more words your clue relates to, the quicker you can win the game, but the harder the associations may be to guess for your partner, so choose your clue word wisely. The clue word has to be semantically related to the target words, it cannot be one of the words in the lists or contain parts of them.} \\
\\ 
\texttt{Always give your single{-}word clue and your comma{-}separated list of related target words in the following format and make your answers as short as possible, never include any other text than is required in this form:} \\
\\ 
\texttt{CLUE: <WORD>} \\
\texttt{TARGETS: <TARGETS>} \\
\\ 
\texttt{Your team words are: polyester, gymnastics, france, birch, hat.} \\
\texttt{Your opponent's team words are: .} \\
\texttt{Distractor words are: ostrich, volleyball, korea, pants, archbishop, acid.} \\
\texttt{The kill words are: coal.} \\
            }
        }
    }
    & \\ \\

    \theutterance \stepcounter{utterance}  
    & \multicolumn{4}{p{0.6\linewidth}}{
        \cellcolor[rgb]{0.9,0.9,0.9}{
            \makecell[{{p{\linewidth}}}]{
                \texttt{\tiny{[P1$\rangle$GM]}}
                \texttt{CLUE: olympics} \\
\texttt{TARGETS: gymnastics, france} \\
            }
        }
    }
    & & \\ \\

    \theutterance \stepcounter{utterance}  
    & & & \multicolumn{2}{p{0.3\linewidth}}{
        \cellcolor[rgb]{0.9,0.9,0.9}{
            \makecell[{{p{\linewidth}}}]{
                \texttt{\tiny{[GM$|$GM]}}
                \texttt{olympics} \\
            }
        }
    }
    & & \\ \\

    \theutterance \stepcounter{utterance}  
    & & & \multicolumn{2}{p{0.3\linewidth}}{
        \cellcolor[rgb]{0.9,0.9,0.9}{
            \makecell[{{p{\linewidth}}}]{
                \texttt{\tiny{[GM$|$GM]}}
                ['gymnastics', 'france']
            }
        }
    }
    & & \\ \\

    \theutterance \stepcounter{utterance}  
    & & & \multicolumn{2}{p{0.3\linewidth}}{
        \cellcolor[rgb]{0.9,0.9,0.9}{
            \makecell[{{p{\linewidth}}}]{
                \texttt{\tiny{[GM$|$GM]}}
                {'word': 'gymnastics', 'assignment': 'team'}
            }
        }
    }
    & & \\ \\

    \theutterance \stepcounter{utterance}  
    & & & \multicolumn{2}{p{0.3\linewidth}}{
        \cellcolor[rgb]{0.9,0.9,0.9}{
            \makecell[{{p{\linewidth}}}]{
                \texttt{\tiny{[GM$|$GM]}}
                {'word': 'france', 'assignment': 'team'}
            }
        }
    }
    & & \\ \\

    \theutterance \stepcounter{utterance}  
    & & & \multicolumn{2}{p{0.3\linewidth}}{
        \cellcolor[rgb]{0.9,0.9,0.9}{
            \makecell[{{p{\linewidth}}}]{
                \texttt{\tiny{[GM$|$GM]}}
                \texttt{opponent has won} \\
            }
        }
    }
    & & \\ \\

    \theutterance \stepcounter{utterance}  
    & \multicolumn{4}{p{0.6\linewidth}}{
        \cellcolor[rgb]{0.95,0.95,0.95}{
            \makecell[{{p{\linewidth}}}]{
                \texttt{\tiny{[P2$\rightarrow$P1]}}
                \texttt{Post{-}episode Reflexion (not sent to GM):} \\
\texttt{New plan: My previous attempt failed because the clue "olympics" was too broad, likely causing my partner to guess "volleyball" (a distractor) or "korea" (an opponent word), which ended our turn or caused a loss. Additionally, I only targeted two words, leaving "polyester," "birch," and "hat" unguessed, which slowed our progress. Next time, I will prioritize high{-}specificity clues that link multiple target words while strictly avoiding semantic overlap with distractors (ostrich, volleyball, korea, pants, archbishop, acid) and opponent words. I will group "polyester" and "hat" under the clue "fabric" (avoiding "pants" by ensuring the context implies clothing material rather than specific garments, or by using "synthetic" if "fabric" is too risky). I will group "gymnastics" and "france" under "host" (referencing France hosting the Olympics) or "rings" (specific to gymnastics, though this might only hit one). A safer multi{-}word strategy is to use "material" for "polyester" and "birch" (as in wood material), and "sport" for "gymnastics" only if "volleyball" is not a risk, or better yet, use "gym" for "gymnastics" and "country" for "france" separately if grouping is too risky. I will ensure every clue covers at least two words if possible, or three, to maximize efficiency, and I will double{-}check that the clue does not trigger "volleyball" (e.g., avoid "sport" if "volleyball" is a distractor) or "korea". Specifically, I will try "synthetic" for "polyester" and "hat" (if the hat is synthetic) or "tree" for "birch" (if "pine" is not a distractor, but "ostrich" is the only bird distractor, so "tree" might be safe for "birch" alone, but I need to group). Actually, "birch" and "polyester" don't group well. "Gymnastics" and "france" are the strongest pair via "olympics" but that failed. I will try "paris" for "france" and "gymnastics" (Paris hosted gymnastics), which is more specific than "olympics" and less likely to trigger "volleyball" or "korea". For the remaining, I will use "fabric" for "polyester" and "hat". I must ensure "fabric" doesn't trigger "pants". If "fabric" is risky, I will use "synthetic" for "polyester" and "headwear" for "hat". The key is to avoid the broad "olymp} \\
            }
        }
    }
    & & \\ \\

\end{supertabular}
}

\end{document}
